	% ============================================================
	% 2. OBJETIVO GENERAL Y ESPECÍFICOS
	% ============================================================
	\section{Objetivo general y específicos (máximo 1 página)}
	
	\subsection{Objetivo general}
	Desarrollar un modelo de aprendizaje profundo basado en la arquitectura YOLO + LSTM y el framework PCB (\textit{Pre Crime Behavior}), optimizado para bajo costo computacional y desplegable en edge devices, para mejorar las métricas de rendimiento en la detección temprana de conductas de \textit{shoplifting} en el entorno retail minorista peruano.
	
	\subsection{Objetivos específicos}

		\begin{enumerate}
			\item \textbf{Construir y estructurar un dataset propio} con diversas secuencias de video de comportamiento de compra y tambien conductas sospechosas en diferente comercios minoristas locales, aplicando diferntes técnicas de preprocesamiento, segmentación de secuencias y data augmentation.
			
			\item \textbf{Diseñar e integrar la arquitectura híbrida YOLO + LSTM}, combinando redes CNN y LSTM para la extracción de características espaciales y el modelado temporal de la secuencia pre-delictiva, consiguiendo tambien optimizar hiperparámetros como la longitud de secuencia y el tamaño de lote.
			
			\item \textbf{Implementar y evaluar el despliegue del modelo sobre edge diveces} de bajo consumo computacional como en tarjetas NVIDIA Jetson o Raspberry Pi, analizando las diferente métricas de eficiencia operativa como la latencia de detección y los fotogramas por seguno (FPS).
			
			\item \textbf{Validar la precisión y rendimiento predictivo del sistema} mediante las diversas métricas estándar como Precisión, Recall, F1-Score y mAP, comparando los resultados obtenidos frente a los detectores espaciales estáticos y arquitecturas pesadas reportadas en la literatura.
		\end{enumerate}
	
